% Source: https://artofproblemsolving.com/community/c3h1376545_can_u_solve_this_easy_task
% Stable identity: aops-topic-1376545-q01
Extend $AP$ beyond $P$ to meet the line $BC$ at $Q$.

\begin{asy}
usepackage("amsmath");
size(7cm);
defaultpen(fontsize(10pt));
pair A=(0,0), D=(8,0), B=(1.5,3), C=(5,3);
pair P=(C+D)/2, Q=(13,3);
draw(A--B--C--D--cycle,linewidth(0.9));
draw(C--Q,gray+dashed); draw(P--Q,gray+dashed);
fill(A--B--P--cycle,paleblue); draw(A--B--P--cycle,blue);
draw(A--B--C--D--cycle,linewidth(0.9));
dot(P); dot(Q);
label("$A$",A,SW); label("$D$",D,SE); label("$B$",B,NW); label("$C$",C,N);
label("$P$",P,1.6*dir(62)); label("$Q$",Q,NE);
label("\mbox{}",(min(currentpicture,user=true)+max(currentpicture,user=true))/2,black); // keeps the text after the figure black when LaTeX typesets the labels inline
\end{asy}

Triangles $PDA$ and $PCQ$ are congruent: $PD=PC$, the vertical angles at $P$ are equal, and $\angle PDA=\angle PCQ$ because $AD\parallel CQ$ (alternate angles). Hence $[PDA]=[PCQ]$ and $AP=PQ$. Therefore
\[
[ABQ]=[ABCP]+[PCQ]=[ABCP]+[PDA]=[ABCD]=60 .
\]
In triangle $ABQ$, $BP$ is a median ($P$ is the midpoint of $AQ$), and a median splits a triangle into two triangles of equal area. So
\[
[ABP]=\dfrac{1}{2}[ABQ]=\boxed{30}.
\]
(If instead $AB$ and $CD$ were the parallel sides, the area of $ABP$ would be $\dfrac{1}{2}\cdot AB\cdot h$, which is not determined by the area of the trapezoid alone.)
